% Fase 1 --- Descrição do produto e modelo de qualidade (responsável: Hugo).
% Pacotes usados: tikz (bibliotecas positioning, fit, backgrounds, arrows.meta,
% shapes.geometric) e url. Versão completa para a GitPage: docs/fase1-produto.md.

\subsection{Descrição do Produto}
\label{subsec:produto}

O objeto da avaliação é a consulta de infrações de trânsito oferecida pelo Detran-DF em dois canais: o Portal de Serviços, apresentado como ``Detran Digital'' \cite{detrandf2026portal}, e o aplicativo ``DETRAN-DF'', versão 2.6.9, distribuído na Google Play e na App Store \cite{detrandf2026app}. O levantamento, feito em 18/10/2026, usou apenas informações públicas: telas do portal, páginas das lojas e os arquivos que o navegador baixa ao abrir o portal; o registro completo está no repositório.

Trata-se de um serviço digital governamental ao cidadão, prestado por um sistema de informação com dois clientes. O portal é uma aplicação web de página única executada no navegador; o aplicativo é um cliente móvel multiplataforma em React Native para Android e iOS. Nenhum dos dois guarda as infrações: ambos usam o mesmo servidor de identidade e a mesma família de APIs do Detran-DF, com serviços distintos para a consulta, que reúnem dados próprios e do Registro Nacional de Infrações de Trânsito (RENAINF). A Tabela~\ref{tab:canais} compara os canais e a Figura~\ref{fig:contexto} mostra o contexto do sistema.

\begin{table}[htbp]
\centering
\caption{Comparação entre os canais da consulta de infrações (levantamento de 18/10/2026)}
\label{tab:canais}
\footnotesize
\begin{tabular}{|p{2.5cm}|p{5.5cm}|p{5.5cm}|}
\hline
\textbf{Aspecto} & \textbf{Portal de Serviços} & \textbf{Aplicativo} \\ \hline
Acesso & \url{portal.detran.df.gov.br}, em navegador de desktop ou celular & Google Play e App Store; Android 8.0 ou iOS 15.1 ou superior \\ \hline
Autenticação & CPF/CNPJ e senha com reCAPTCHA, ou conta gov.br & Mesma conta: CPF/CNPJ e senha, gov.br ou biometria, sem reCAPTCHA \\ \hline
Abrangência da consulta & Só veículos vinculados ao CPF autenticado; opção separada para veículo de outra UF & Igual ao portal \\ \hline
Dados por infração & 31 rótulos previstos na tela de detalhes, entre eles auto, enquadramento, local, data, hora, órgão autuador, natureza, pontuação, valores, prazos e situação & Os mesmos rótulos, exceto UF, mais a data de opção pelo desconto do SNE (a confirmar com veículo real) \\ \hline
\end{tabular}
\par\smallskip
\raggedright\scriptsize Fonte: elaborada pela equipe a partir de \cite{detrandf2026portal} e \cite{detrandf2026app}.
\end{table}

\begin{figure}[htbp]
\centering
% Diagrama de contexto da consulta de infrações (Figura fig:contexto).
% Fonte editável em TikZ. Requer no preâmbulo:
%   \usepackage{tikz}
%   \usetikzlibrary{positioning,fit,backgrounds,arrows.meta}
% Versão Mermaid equivalente: docs/fase1-produto.md, Seção 2.3.
\begin{tikzpicture}[
    font=\scriptsize,
    node distance=4mm and 8mm,
    caixa/.style={draw, rounded corners=2pt, align=center, minimum height=8mm, text width=31mm, fill=white},
    base/.style={draw, cylinder, shape border rotate=90, aspect=0.15, align=center, minimum height=9mm, text width=24mm, fill=white},
    grupo/.style={draw, dashed, rounded corners=3pt, inner sep=2.5mm},
    rotulo/.style={font=\scriptsize\itshape, anchor=south west, inner sep=1pt},
    seta/.style={-{Stealth[length=1.6mm]}},
  ]
  % Usuário
  \node[caixa, text width=22mm] (cidadao) {Cidadão / condutor\\(proprietário)};

  % Canais avaliados
  \node[caixa, right=of cidadao, yshift=7mm] (portal) {Portal de Serviços\\(SPA no navegador)};
  \node[caixa, below=of portal] (app) {Aplicativo DETRAN-DF\\(React Native)};

  % Detran-DF
  \node[caixa, right=of portal, yshift=4mm] (idp) {Servidor de identidade\\(OpenID Connect)};
  \node[caixa, below=of idp] (api) {Serviços de backend\\(API HTTPS)};
  \node[base, below=of api] (base) {Base de infrações\\do Detran-DF};

  % Externos
  \node[caixa, right=of idp, text width=21mm] (govbr) {gov.br\\(login federado)};
  \node[base, right=of api, yshift=-3mm] (renainf) {RENAINF\\(SENATRAN)};
  \node[caixa, above=6mm of portal, text width=31mm] (terceiros) {reCAPTCHA, CDN, Firebase\\e monitoramento};

  \begin{scope}[on background layer]
    \node[grupo, fit=(portal)(app)] (gcanais) {};
    \node[grupo, fit=(idp)(api)(base), fill=black!4] (gdetran) {};
    \node[grupo, fit=(govbr)(renainf)] (gext) {};
  \end{scope}
  \node[rotulo, anchor=north west] at (gcanais.south west) {Canais avaliados};
  \node[rotulo] at (gdetran.north west) {Detran-DF (fora do alcance)};
  \node[rotulo] at (gext.north west) {Externos};

  \draw[seta] (cidadao) -- (portal);
  \draw[seta] (cidadao) -- (app);
  % Os dois canais usam os mesmos serviços: as setas saem do grupo.
  \draw[seta] (gcanais.east |- idp) -- (idp);
  \draw[seta] (gcanais.east |- api) -- (api);
  \draw[seta] (gcanais.north -| terceiros) -- (terceiros);
  \draw[seta] (idp) -- (govbr);
  \draw[seta] (api) -- (base);
  \draw[seta] (api) -- (renainf);
\end{tikzpicture}

\caption{Diagrama de contexto da consulta de infrações.}
\label{fig:contexto}
\par\smallskip
\scriptsize Fonte: elaborada pela equipe.
\end{figure}

Por ser caixa-preta, a avaliação alcança apenas o que chega ao avaliador: campos exibidos, passos até o resultado, tempo percebido, dados transferidos e conclusão da consulta em cada contexto de uso. Ficam fora o tempo de processamento no servidor, a capacidade da infraestrutura, a disponibilidade ao longo do tempo e a correção absoluta dos dados, sem acesso à base oficial. O produto impõe ainda três condições ao plano: só o proprietário que consentiu pode fazer login, pois a consulta lista apenas veículos do CPF autenticado; o login difere entre canais (reCAPTCHA no portal, biometria no aplicativo), então a medição de tempo começa após a autenticação; e o perfilador do Android Studio não funciona com aplicativos de loja, o que restringe a utilização de recursos ao volume de dados transferidos.

\subsection{Modelo de Qualidade Adaptado}
\label{subsec:modelo-qualidade}

O modelo parte de duas fontes da família SQuaRE. Adequação funcional e eficiência de desempenho vêm do modelo de qualidade do produto da ISO/IEC 25010:2023 \cite{iso25010_2023}. A cobertura de contexto vem do modelo de qualidade em uso da ISO/IEC 25010:2011 \cite{iso25010_2011}, com medidas na ISO/IEC 25022:2016 \cite{iso25022_2016}. O uso da edição anterior é deliberado: a ISO/IEC 25019:2023, que substituiu esse modelo, removeu a cobertura de contexto \cite{iso25019_2023}, mas a comparação entre canais depende justamente dos contextos em que cada um é usado, e as medidas correspondentes continuam vigentes na ISO/IEC 25022:2016.

A Tabela~\ref{tab:modelo-qualidade} registra a decisão sobre cada subcaracterística das três características escolhidas e a parte interessada que a motiva. As demais características da ISO/IEC 25010:2023 ficam fora: capacidade de interação, por corresponder à antiga usabilidade; compatibilidade, porque a troca de dados com gov.br e RENAINF ocorre no servidor; confiabilidade, porque disponibilidade exige monitoramento por semanas; segurança, que exigiria autorização formal do órgão; e manutenibilidade, flexibilidade e segurança operacional, que dependem de código-fonte ou não se aplicam a uma consulta de dados. Quatro adaptações distinguem o modelo do original. A completude de contexto é medida só pela conclusão da consulta e pelo tempo, sem satisfação, que pertence à usabilidade e está fora do escopo. A correção funcional é julgada pela consistência entre os canais e com a notificação recebida pelo proprietário. A apropriação funcional segue o exemplo da própria norma, que fala em passos necessários e suficientes para concluir a tarefa, e é medida pela contagem objetiva de passos. A utilização de recursos se restringe à rede. Para a completude funcional, propõe-se uma lista de referência externa aos dois canais, derivada do conteúdo obrigatório do auto de infração (art.~280 do Código de Trânsito Brasileiro) e das informações de natureza, pontuação e valor (arts.~258 e 259) \cite{brasil1997ctb}. A Figura~\ref{fig:modelo-qualidade} apresenta a hierarquia resultante.

\begin{table}[htbp]
\centering
\caption{Seleção das subcaracterísticas do modelo de qualidade adaptado}
\label{tab:modelo-qualidade}
\footnotesize
\begin{tabular}{|p{2.9cm}|p{1.6cm}|p{9.1cm}|}
\hline
\textbf{Subcaracterística} & \textbf{Decisão} & \textbf{Justificativa} \\ \hline
\multicolumn{3}{|l|}{\textit{Adequação funcional (ISO/IEC 25010:2023)}} \\ \hline
Completude funcional & Incluída & O cidadão precisa de todos os dados da autuação para pagar, recorrer ou indicar condutor \\ \hline
Correção funcional & Adaptada & Divergência entre canais leva a pagamento errado ou perda de prazo; medida como consistência entre canais \\ \hline
Apropriação funcional & Adaptada & Passos desnecessários atrasam o cidadão; medida pela contagem de passos, sem opinião \\ \hline
\multicolumn{3}{|l|}{\textit{Eficiência de desempenho (ISO/IEC 25010:2023)}} \\ \hline
Comportamento temporal & Incluída & A equipe técnica do Detran-DF precisa comparar tempos entre canais e redes \\ \hline
Utilização de recursos & Adaptada & Só dados transferidos, que pesam no plano de dados do cidadão e são medidos igualmente nos dois canais \\ \hline
Capacidade & Excluída & Exige teste de carga sobre a infraestrutura do órgão \\ \hline
\multicolumn{3}{|l|}{\textit{Cobertura de contexto (ISO/IEC 25010:2011, qualidade em uso)}} \\ \hline
Completude de contexto & Adaptada & O requisitante precisa saber se a consulta funciona em todos os dispositivos e redes previstos; sem satisfação \\ \hline
Flexibilidade & Incluída & Cidadãos com rede degradada ou sistema antigo, fora dos contextos previstos \\ \hline
\end{tabular}
\par\smallskip
\raggedright\scriptsize Fonte: elaborada pela equipe com base em \cite{iso25010_2023} e \cite{iso25010_2011}.
\end{table}

\begin{figure}[htbp]
\centering
% Modelo de qualidade adaptado (Figura fig:modelo-qualidade).
% Fonte editável em TikZ. Requer no preâmbulo:
%   \usepackage{tikz}
%   \usetikzlibrary{positioning,fit,backgrounds,arrows.meta}
% Versão Mermaid equivalente: docs/fase1-produto.md, Seção 3.4.
\begin{tikzpicture}[
    font=\scriptsize,
    raiz/.style={draw, thick, rounded corners=2pt, align=center, text width=60mm, minimum height=8mm, fill=white},
    carac/.style={draw, rounded corners=2pt, align=center, text width=30mm, minimum height=7mm, fill=white, font=\scriptsize\bfseries},
    sub/.style={draw, align=center, text width=26mm, minimum height=6mm, fill=white},
    grupo/.style={draw, rounded corners=3pt, inner sep=2.5mm},
    rotulo/.style={font=\scriptsize\itshape, anchor=north, inner sep=2pt, align=center},
    liga/.style={-{Stealth[length=1.4mm]}},
  ]
  \node[raiz] (raiz) {Qualidade da consulta de infrações\\Portal de Serviços $\times$ aplicativo};

  % Características
  \node[carac, below left=12mm and 18mm of raiz.south] (af) {Adequação funcional};
  \node[carac, right=6mm of af] (ed) {Eficiência de desempenho};
  \node[carac, right=12mm of ed] (cc) {Cobertura de contexto};

  % Subcaracterísticas
  \node[sub, below=5mm of af, xshift=3mm] (af1) {Completude funcional};
  \node[sub, below=2mm of af1] (af2) {Correção funcional*};
  \node[sub, below=2mm of af2] (af3) {Apropriação funcional*};
  \node[sub, below=5mm of ed, xshift=3mm] (ed1) {Comportamento temporal};
  \node[sub, below=2mm of ed1] (ed2) {Utilização de recursos*};
  \node[sub, below=5mm of cc, xshift=3mm] (cc1) {Completude de contexto*};
  \node[sub, below=2mm of cc1] (cc2) {Flexibilidade};

  \begin{scope}[on background layer]
    \node[grupo, fill=black!3, fit=(af)(ed)(af3)] (gprod) {};
    \node[grupo, dashed, fit=(cc)(cc2)] (guso) {};
  \end{scope}
  \node[rotulo] at (gprod.south) {Qualidade do produto\\ISO/IEC 25010:2023};
  \node[rotulo] at (guso.south) {Qualidade em uso\\ISO/IEC 25010:2011\\(medidas: ISO/IEC 25022:2016)};

  \foreach \c in {af,ed,cc} \draw[liga] (raiz.south) -- (\c.north);
  \foreach \s in {af1,af2,af3} \draw[liga] ([xshift=-13mm]af.south) |- (\s.west);
  \foreach \s in {ed1,ed2} \draw[liga] ([xshift=-13mm]ed.south) |- (\s.west);
  \foreach \s in {cc1,cc2} \draw[liga] ([xshift=-13mm]cc.south) |- (\s.west);
\end{tikzpicture}

\caption{Modelo de qualidade adaptado. O asterisco indica subcaracterística adaptada.}
\label{fig:modelo-qualidade}
\par\smallskip
\scriptsize Fonte: elaborada pela equipe com base em \cite{iso25010_2023} e \cite{iso25010_2011}.
\end{figure}
